\chapter{نتیجه‌گیری}

در این گزارش، مفاهیم بنیادی و عملیِ دو حوزهٔ مرتبط با مدل‌های مولد عمیق تشریح شد: نخست، مدل‌های گرافی احتمالاتی (PGMs) که ابزارهای نظری و عملی برای نمایش و استنتاج در مسائل با عدم قطعیت ارائه می‌دهند؛ و دوم، اتوانکدرهای واریاسیونال (VAE) که چارچوبی برای یادگیری توزیع‌های مولد بر پایهٔ تقریب واریانسی فراهم می‌آورند.

\section{خلاصهٔ نتایج و پیامدها}
\begin{itemize}
    \item PGMs نمایش ساختاریِ وابستگی‌ها را مقدور می‌سازند و با قواعد استقلال شرطی امکان ساده‌سازیِ محاسبات پیچیده را فراهم می‌کنند. ابزارهای استنتاج دقیق و تقریب‌پذیر بسته به ساختار گراف و منابع محاسباتی انتخاب می‌شوند.
    \item VAEها با استفاده از ELBO و باز-پارامتری‌سازی، آموزشِ مؤثرِ مدل‌های مولد را تسهیل می‌کنند؛ تنظیمِ ضریبِ KL (مثلاً \beta) و استفاده از جریان‌ها می‌تواند تعادل بین بازسازی و disentanglement را تغییر دهد.
    \item در عمل، انتخاب معماری و هایپرپارامترها نیازمند آزمون‌های تجربی و توجه به معیارهای کمی (MSE, ELBO, MIG) و کیفی (بازسازی بصری، تراورس) است.
\end{itemize}

\section{محدودیت‌ها و ملاحظات اخلاقی}
\begin{itemize}
    \item از منظر محاسباتی، مدل‌های پیچیده نیاز به منابع محاسباتی و زمان آموزش بالایی دارند. برآورد و مدیریت عدم قطعیت در پیش‌بینی‌ها برای کاربردهای حسّاس ضروری است.
    \item تولید محتوای مصنوعی می‌تواند مسائل اخلاقی و حریم خصوصی ایجاد کند؛ در کاربردهای عملی نیاز به کنترل کیفیت، جلوگیری از سوءاستفاده و توجه به ملاحظات داده‌ای است.
\end{itemize}

\section{پیشنهادات پژوهشی و کاربردی}
\begin{itemize}
    \item ارزیابی سیستماتیک روی مجموعهٔ داده‌های مختلف برای سنجش قابلیت تعمیم‌پذیری روش‌ها.
    \item توسعهٔ روش‌های ترکیبی (مثلاً VAE+Flow یا VAE+GAN) برای بهبود کیفیت نمونه‌ها.
    \item ایجاد چارچوب‌های اتوماتیک برای انتخاب \beta و ساختار فضای پنهان بر اساس هدف مسئله.
\end{itemize}

\section{جمع‌بندی نهایی}
این گزارش چارچوب نظری و عملی مناسبی برای ادامهٔ کار در حوزهٔ مدل‌های مولد فراهم می‌کند. خواننده می‌تواند با پیاده‌سازی تجربی و ارزیابی کمیِ پیشنهادات مطرح‌شده، به نتایج علمی و کاربردی قابل توجهی دست یابد.
